% subsection: 等差型 \, $a_{n+1}=a_n+d$

  \subsection{\texorpdfstring{等差型 \, $a_{n+1}=a_n+d$}{等差型}}\label{節：等差型}
  \begin{bookpoint}
    \textbf{差を追う}\par
    一回ごとに足す量が一定なら,初項に増分を足し戻す.項の番号と更新の回数の違いを確かめよう.
  \end{bookpoint}
    まずはもっともシンプルな\bookterm{recurrence}{漸化式}であろう,等差型について説明していく.
    この\bookterm{recurrence}{漸化式}は$a_1$に対して$d$がどんどんと足されていく\bookterm{recurrence}{漸化式}である.
    \bookterm{general}{一般項}は簡単な代入計算によって容易に示せる.
    実際にやってみると
    \begin{align*}
      &a_{n+1}=a_n+d\\
      &a_{n+1}=(a_{n-1}+d)+d\\
      &a_{n+1}=\{(a_{n-2}+d)+d\}+d\\
      &\qquad \vdots\\
      &a_{n+1}=a_1+d+d+\cdots+d\\
      &a_{n+1}=a_1+nd\\
      &\therefore \quad a_n=a_1+(n-1)d
    \end{align*}
    となり\bookterm{general}{一般項}が示される.
    1つ前の項をどんどん代入していくのだ.
    数式の3行目から5行目周辺の式変形が少し心配かもしれないが,\,$d$の個数と$a$の添字部分の和は常に$n+1$となることから考えれば分かるだろう.
    実際,\,$a_{n+1}$に至るまでに何度$d$が足されたかを考えているわけだから,心配であれば実際に適当な$n$で何度か実験してみるのもいい手だろう.

    実際に例題を与えよう.

    \noindent \bookblocklink{example-arithmetic-1}{problem}{answer}{【例題】}\par\nobreak
    \begin{enumerate}[label=(\arabic*),parsep=3pt,format=\bookitemlink{example-arithmetic-1}{problem}{answer}]
      \item $\displaystyle a_1=1,\, a_{n+1}=a_n+1$
      \item $\displaystyle a_1=\frac{1}{2},\, a_{n+1}=a_n+\frac{1}{4}$
      \item $\displaystyle a_1=-\sqrt{3},\, a_{n+1}=a_n-\sqrt{2}$
    \end{enumerate}

    \noindent\bookblocklink{example-arithmetic-1}{hint}{problem}{【方針】}\par\nobreak
    \begin{enumerate}[label=(\arabic*),parsep=3pt,format=\bookitemlink{example-arithmetic-1}{hint}{problem}]
      \item 等差型の\bookterm{general}{一般項}
            $\displaystyle a_n=a_1+(n-1)d$
            を用います.
      \item 公差は$\displaystyle d=\frac14$です.
            等差型の\bookterm{general}{一般項}に代入します.
      \item 公差は$\displaystyle d=-\sqrt2$です.
            負の公差の場合も同じ公式を使います.
    \end{enumerate}

    \noindent\bookblocklink{example-arithmetic-1}{answer}{problem}{【解答】}\par\nobreak
    \begin{enumerate}[label=(\arabic*),parsep=3pt,format=\bookitemlink{example-arithmetic-1}{answer}{problem}]
      \item $\displaystyle a_n=1+(n-1)=n$
      \item $\displaystyle a_n=\frac12+\frac{n-1}{4}
            =\frac{n+1}{4}$
      \item $\displaystyle a_n=-\sqrt3-(n-1)\sqrt2$
    \end{enumerate}

    \noindent\bookblocklink{example-arithmetic-1}{solution}{problem}{【解法】}\par\nobreak
    解法略.
